\section*{Bloque C. Casos con implicancias regulatorias y sector saneamiento}

\noindent\textit{Preguntas 66 a 90. Evalúan aplicaciones al marco tarifario y financiero de SUNASS y saneamiento peruano: WACC regulatorio, Base de Activos Regulatorios (BAR), Valor Nuevo de Reemplazo (VNR), evaluación social en Invierte.pe (Tasa Social de Descuento del 8\%), contratos APP en PTAR, descalce cambiario y sostenibilidad de EPS.}

\q{En la metodología regulatoria de SUNASS para la determinación del Costo Promedio Ponderado de Capital (WACC regulatorio) aplicado a las tarifas quinquenales de agua:}
\begin{opts}
  \item se asume que el costo del capital propio es igual a la inflación del INEI.
  \item se fija una estructura de capital regulatoria de referencia normativa (típicamente 50\% deuda y 50\% capital propio), incentivando a la EPS a optimizar su financiamiento sin trasladar ineficiencias de endeudamiento a los usuarios.
  \item se utiliza la estructura contable de deuda histórica de cada municipio.
  \item se prohíbe el uso de endeudamiento bancario.
\end{opts}

\q{Para estimar la Tasa Libre de Riesgo ($R_f$) en el CAPM del WACC de las EPS, SUNASS toma como parámetro de referencia internacional:}
\begin{opts}
  \item la tasa de inflación de Lima Metropolitana.
  \item la tasa de las cuentas de ahorro del Banco de la Nación.
  \item el rendimiento promedio de los Bonos del Tesoro de los Estados Unidos (US Treasury Bonds) a 10 o 30 años, por ser el activo con menor riesgo de crédito a nivel global.
  \item el rendimiento de los bonos corporativos de empresas mineras.
\end{opts}

\q{El spread de Riesgo País incorporado por SUNASS en el costo de oportunidad del capital propio de las EPS se mide usualmente a través de:}
\begin{opts}
  \item la tasa de desempleo juvenil de Lima.
  \item el arancel promedio a la importación de bienes de capital.
  \item el índice de precios al consumidor de Sedapal.
  \item el indicador EMBI+ Perú (Emerging Markets Bond Index) elaborado por J.P. Morgan, que refleja el diferencial de rendimiento entre los bonos soberanos peruanos y los bonos del tesoro norteamericano.
\end{opts}

\q{La Base de Activos Regulatorios (BAR / RAB) sobre la cual SUNASS aplica la tasa WACC para remunerar el capital en las tarifas de las EPS comprende:}
\begin{opts}
  \item el valor económico de los activos fijos netos en operación directamente afectos y utilizados en la prestación eficiente de los servicios de agua potable y alcantarillado.
  \item el total de activos intangibles y patentes de software de la empresa.
  \item el monto total de las cuentas por cobrar incobrables de usuarios morosos.
  \item las inversiones financieras especulativas en fondos mutuos.
\end{opts}

\q{Para valorar la infraestructura de redes subterráneas y plantas antiguas de agua, SUNASS y la teoría regulatoria aplican la metodología del Valor Nuevo de Reemplazo (VNR), la cual valora los activos según:}
\begin{opts}
  \item el precio de venta de las acciones en bolsa.
  \item el costo corriente que demandaría construir y poner en operación hoy un activo moderno equivalente con tecnología eficiente y capacidad adecuada, deduciendo la depreciación acumulada por estado de conservación.
  \item el valor de remate de la chatarra metálica.
  \item su costo histórico de adquisición contable registrado hace 40 años.
\end{opts}

\q{Aunque las 50 EPS del Perú son empresas públicas municipales que no buscan maximizar beneficios como una corporación privada, el regulador SUNASS les reconoce obligatoriamente una tasa WACC positiva en las tarifas para:}
\begin{opts}
  \item trasladar utilidades al Tesoro Público central.
  \item permitir el pago de bonos de productividad ilimitados a los funcionarios municipales.
  \item asegurar la sostenibilidad financiera a largo plazo, la reposición oportuna de activos depreciados y la capacidad de apalancar crédito para financiar la expansión de cobertura.
  \item financiar campañas de publicidad electoral en las provincias.
\end{opts}

\q{En el Sistema Nacional de Programación Multianual y Gestión de Inversiones (Invierte.pe), la evaluación de proyectos públicos de saneamiento se realiza bajo el enfoque de:}
\begin{opts}
  \item análisis de opciones financieras de derivados OTC.
  \item evaluación financiera privada exclusiva con precios de mercado.
  \item aprobación automática sin formulación de flujos de caja.
  \item evaluación social de proyectos, corrigiendo los precios de mercado mediante Factores de Corrección Social para reflejar el verdadero costo de oportunidad para la sociedad peruana.
\end{opts}

\q{La Tasa Social de Descuento (TSD) fijada por el Ministerio de Economía y Finanzas (MEF) en los anexos metodológicos de Invierte.pe para la evaluación de proyectos de saneamiento es actualmente de:}
\begin{opts}
  \item $8.0\%$ anual en términos reales
  \item $15.0\%$ anual
  \item $3.0\%$ anual
  \item $12.0\%$ anual
\end{opts}

\q{En la evaluación social de un proyecto de zanjas y reservorios en una zona andina deprimida, el Factor de Corrección Social de la Mano de Obra No Calificada (FC = 0.41 a 0.60 según región) es menor a 1 porque:}
\begin{opts}
  \item el Estado exonera del pago de remuneraciones a los campesinos.
  \item el salario de mercado supera el verdadero costo de oportunidad social del trabajador debido al desempleo o subempleo rural, por lo que el costo social para el país es menor al gasto contable de la planilla.
  \item los trabajadores rurales son más productivos que los ingenieros.
  \item la mano de obra se financia con donaciones externas.
\end{opts}

\q{Cuando un proyecto de agua potable rural en una comunidad nativa remota genera beneficios de reducción de parasitosis difíciles de cuantificar en soles, la metodología de evaluación de Invierte.pe que debe aplicarse es:}
\begin{opts}
  \item la TIR de mercado del proyecto.
  \item el Análisis Costo-Beneficio monetario con VAN social.
  \item el Análisis Costo-Eficacia o Costo-Efectividad, seleccionando la alternativa técnica que logre la meta de provisión y calidad al menor Costo Promedio por beneficiario o por metro cúbico.
  \item el remate público de la infraestructura.
\end{opts}

\q{Los bancos multilaterales (BID, Banco Mundial, KfW) que otorgan créditos de inversión a las EPS exigen el cumplimiento de un 'Ratio de Cobertura del Servicio de la Deuda' (DSCR), el cual se define como:}
\begin{opts}
  \item $DSCR = \frac{Pasivo\ Total}{Ventas\ Anuales}$
  \item $DSCR = \frac{Activo\ Fijo}{Patrimonio}$
  \item $DSCR = \frac{Utilidad\ Neta}{EBITDA}$
  \item $DSCR = \frac{Flujo\ de\ Caja\ Operativo\ Neto}{Servicio\ de\ la\ Deuda\ (Principal + Intereses)}$, exigiendo usualmente que sea $\ge 1.20$ veces para garantizar el repago.
\end{opts}

\q{El 'Periodo Promedio de Cobranza' de una EPS que supera los 120 días comerciales refleja un grave problema financiero ocasionado frecuentemente por:}
\begin{opts}
  \item la elevada morosidad y retrasos de pago sistemáticos de instituciones públicas (ministerios, comisarías, colegios estatales y hospitales) e industrias, inmovilizando capital de trabajo y forzando iliquidez.
  \item que los usuarios pagan sus facturas con un año de anticipación.
  \item un exceso de inversión en micromedidores ultrasónicos.
  \item la indexación tarifaria mensual automática.
\end{opts}

\q{La 'Depreciación Regulatoria' incluida por SUNASS en la estructura de costos de la tarifa quinquenal tiene como fin financiero principal:}
\begin{opts}
  \item financiar el subsidio cruzado de los usuarios industriales.
  \item constituir el Fondo de Inversión y Reposición de Activos de la EPS, garantizando que cuando una tubería o bomba culmine su vida útil regulatoria existan recursos monetarios tangibles para reemplazarla.
  \item distribuir utilidades extraordinarias al alcalde provincial.
  \item generar un escudo tributario para no pagar impuesto a la renta.
\end{opts}

\q{En los contratos de concesión de Asociaciones Público-Privadas (APP) para Plantas de Tratamiento de Aguas Residuales (como la PTAR Titicaca), la retribución al concesionario privado se estructura típicamente mediante dos componentes:}
\begin{opts}
  \item el cobro de una comisión por corte y reconexión en toda la región.
  \item la entrega de terrenos agrícolas en propiedad.
  \item la Remuneración por Inversiones (RPI), que amortiza el CAPEX de la planta, y la Remuneración por Mantenimiento y Operación (RPMO), que cubre el OPEX sujeto al cumplimiento de estándares de calidad del efluente.
  \item un arancel a la exportación y un crédito fiscal de IGV.
\end{opts}

\q{En el financiamiento de grandes plantas desalinizadoras o PTAR bajo la modalidad de 'Project Finance' (Financiamiento de Proyecto sin recurso):}
\begin{opts}
  \item los patrocinadores garantizan el 100\% de la deuda con su patrimonio corporativo personal.
  \item el Ministerio de Vivienda asume la totalidad del riesgo geológico.
  \item la tasa de interés del préstamo es subsidiada a tasa cero por la SBS.
  \item el repago de la deuda bancaria depende exclusivamente de los flujos de caja operativos generados por el propio proyecto y los activos de la sociedad de propósito especial (SPE) actúan como única garantía colateral.
\end{opts}

\q{Si una EPS municipal registra pérdidas contables acumuladas sistemáticas y no paga impuesto a la renta corporativo:}
\begin{opts}
  \item el beneficio del escudo fiscal de la deuda ($T \times K_d$) se pierde o se reduce drásticamente, haciendo que el financiamiento mediante endeudamiento sea financieramente menos ventajoso.
  \item los bancos le otorgan préstamos a tasa libre de riesgo.
  \item su WACC regulatorio se vuelve automáticamente negativo.
  \item la empresa debe disolverse obligatoriamente al término del año fiscal.
\end{opts}

\q{Los fondos de reserva para Mecanismos de Retribución por Servicios Ecosistémicos (MERESE) recaudados obligatoriamente en la tarifa de agua por SUNASS deben ser colocados por las EPS en:}
\begin{opts}
  \item el presupuesto de festividades del municipio provincial.
  \item cuentas bancarias intangibles y separadas en entidades del sistema financiero supervisadas por la SBS, utilizables exclusivamente para la ejecución de proyectos de conservación de cuencas aprobados.
  \item la cuenta corriente general para pago de combustibles de camionetas de la empresa.
  \item fondos de cobertura de alto riesgo en criptomonedas.
\end{opts}

\q{Una EPS que contrata un crédito financiero internacional en dólares estadounidenses para construir una planta potabilizadora mientras recauda sus ingresos tarifarios íntegramente en soles peruanos asume un severo riesgo de:}
\begin{opts}
  \item disminución del consumo de agua de los hogares.
  \item arbitraje de tasas de interés de Paridad Descubierta.
  \item descalce cambiario (Currency Mismatch), donde una depreciación del sol incrementa automáticamente la carga del servicio de la deuda en soles sin un aumento correlativo en los ingresos regulados.
  \item pérdida de la titularidad legal sobre las tuberías.
\end{opts}

\q{Al diseñar una Planta de Tratamiento modular donde la Etapa 1 se construye hoy para 500 litros/segundo y existe la 'Opción Real de Expansión' para duplicar a 1000 litros/segundo dentro de 8 años si la población crece:}
\begin{opts}
  \item SUNASS sanciona a la empresa por no construir la planta completa el primer día.
  \item la tasa de descuento social se eleva al 20\%.
  \item el costo total del proyecto se duplica innecesariamente en el año 0.
  \item se reduce el riesgo de sobreinversión en capacidad ociosa no remunerada y se preserva flexibilidad estratégica, maximizando el valor presente ajustado del proyecto.
\end{opts}

\q{Cuando el Organismo Técnico de la Administración de los Servicios de Saneamiento (OTASS) interviene una EPS y la incorpora al Régimen de Apoyo Transitorio (RAT), uno de los indicadores financieros determinantes es:}
\begin{opts}
  \item la causal de insolvencia financiera y patrimonio neto negativo, donde las pérdidas acumuladas han consumido la totalidad del capital social impidiendo la operación autónoma.
  \item el cumplimiento del 100\% de las metas del plan maestro.
  \item el exceso de liquidez corriente en moneda extranjera.
  \item la privatización de las fuentes de captación.
\end{opts}

\q{En el mecanismo de financiamiento de 'Obras por Impuestos' (OxI - Ley N° 29230) aplicado a obras de saneamiento rural:}
\begin{opts}
  \item las obras se ejecutan sin expediente técnico ni supervisión.
  \item una empresa privada financia y ejecuta la obra pública con cargo a deducir el monto invertido contra su pago a cuenta del Impuesto a la Renta de tercera categoría mediante CIPRL.
  \item el municipio exonera del pago de arbitrios a toda la población.
  \item la EPS emite bonos de titulización con garantía del Banco de la Nación.
\end{opts}

\q{La omisión del registro de pasivos ambientales y contingencias judiciales por vertimientos no autorizados en el Balance General de una EPS genera:}
\begin{opts}
  \item la anulación automática de los costos de financiamiento bancario.
  \item la duplicación del margen bruto de ventas.
  \item una sobrestimación del Patrimonio Neto contable de la empresa y una imagen financiera artificialmente saneada que oculta riesgos latentes de multas de la OEFA y SUNASS.
  \item un incremento inmediato de la recaudación tarifaria.
\end{opts}

\q{Al comparar el reemplazo de micromedidores mecánicos convencionales (costo inicial bajo, vida útil de 5 años, submedición por desgaste rápido) frente a medidores ultrasónicos estáticos (costo inicial mayor, vida útil de 12 años, precisión sostenida y telemetría):}
\begin{opts}
  \item los medidores ultrasónicos no pueden ser depreciados tributariamente.
  \item la teoría de las finanzas exige no medir el consumo de agua.
  \item se debe elegir siempre el medidor mecánico por tener menor costo en $t=0$.
  \item el criterio de evaluación financiera idóneo es el Costo Anual Equivalente (CAE), incorporando tanto los costos de reposición como los ingresos adicionales recuperados por reducción de subfacturación.
\end{opts}

\q{El periodo regulatorio tarifario fijado por SUNASS tiene una vigencia estándar de cinco (5) años (quinquenio regulatorio). Esta duración temporal se fundamenta financieramente en:}
\begin{opts}
  \item el ciclo de maduración y ejecución de las inversiones de capital (CAPEX) y la necesidad de otorgar un marco de estabilidad y certidumbre para la recuperación del WACC antes de transferir las ganancias de productividad a los usuarios.
  \item la rotación anual del directorio de la EPS.
  \item el periodo de mandato de los regidores municipales.
  \item la frecuencia del censo nacional del INEI.
\end{opts}

\q{En un contrato de concesión de Asociación Público-Privada (APP) de saneamiento a 25 años, una cláusula de 'Restablecimiento del Equilibrio Económico-Financiero' protege al inversionista privado y al Estado ante:}
\begin{opts}
  \item cambios ordinarios en los precios de mercado que forman parte del riesgo comercial propio del negocio.
  \item eventos sobrevinientes y extraordinarios no atribuibles a las partes, tales como modificaciones normativas regulatorias imprevistas (cambio en la ley), catástrofes naturales de fuerza mayor o variaciones fiscales drásticas.
  \item la decisión unilateral del concesionario de suspender el tratamiento de aguas.
  \item la reducción voluntaria de la jornada laboral de los obreros de la planta.
\end{opts}

